\section{Related Work}
\label{sec:related-work}

\para{Induction heads and in-context learning.} The modern induction-head literature treats these heads as a concrete circuit for copy-based pattern completion. Olsson et al.\ identify induction heads as a likely mechanistic driver of in-context learning and show strong causal evidence in small attention-only models \cite{olsson2022induction}. Crosbie and Shutova extend that claim to larger modern LLMs, where ablating a small fraction of induction heads sharply reduces few-shot gains \cite{crosbie2024induction}. Ren et al.\ broaden the picture further by identifying semantic induction heads that recover structured relations rather than only literal token matches \cite{ren2024semantic}. We build on this line of work by using empirical induction-head scoring to choose ablation targets in \gptsmall, but our target behavior is different: we focus on unwanted copying pressure in prompts that explicitly ask for a diffuse answer.

\para{Induction heads and repetition pathologies.} Wang et al.\ provide the closest prior hypothesis to ours. They argue that induction-head dominance mechanistically explains repetition curse in instruction-tuned LLMs and propose a softer descaling intervention rather than hard ablation \cite{wang2025toxicity}. Their results suggest that induction heads can hurt output diversity in repetitive settings. Our work asks a narrower but complementary question: whether induction-like heads also push probability mass toward a repeated distractor when the task is to choose at random from a constrained set. Unlike Wang et al., we pair this test with random-head controls and a utility check, which makes the trade-off visible.

\para{Memorization, extraction, and distributed leakage.} A separate literature argues against overly local explanations for leakage. Deduplicated training data greatly reduces memorized emissions while improving efficiency, which shows that leakage can be amplified by data duplication rather than by one specific circuit \cite{lee2022deduplicating}. Strong extraction baselines further show that leakage depends on prompting and search strategy, not just on the presence of a memorized string \cite{yu2023bagoftricks}. Ippolito et al.\ show that blocking exact verbatim memorization does not eliminate leakage under modified prompts \cite{ippolito2022falseprivacy}. Huang et al.\ go further and argue that verbatim memorization is intertwined with general language modeling capability, making it hard to isolate without utility loss \cite{huang2024memorization}. Our results align more closely with this distributed view than with a single-circuit story.

\para{Diffuse distributions and random output failures.} Zhang et al.\ show that instruction-tuned LLMs are systematically poor at producing diffuse distributions when prompts ask for random answers \cite{zhang2024diffuse}. That work motivates our random-choice leakage task: if models already struggle to be diffuse, repeated prompt context may further bias them toward copying salient items. Unlike Zhang et al., we do not train a mitigation. Instead, we use mechanistic interventions to ask whether induction-like heads are a major cause of that failure mode.

\para{Positioning.} Our study sits at the intersection of these themes. We combine empirical induction-head detection, causal ablation, a random-choice leakage metric, and a utility control in one real model. The result is a more constrained claim than some mechanistic accounts suggest: induction-like heads matter, but their effect on leakage-like copying is heterogeneous and not cleanly separable from broader model behavior.
